% Sección F · Cobertura de datos. La bitácora y la declaración de IA van en anexos,
% como los pide el enunciado.

\section{Cobertura de datos}

\paragraph{Qué quedó fuera.} De las 13.983 ventas registradas en el ERP quedan fuera \textbf{221, un 1,6\,\%}: 138 duplicados, que sí son un problema de calidad, y 83 facturas de 2026, que quedan en \texttt{fact\_venta} pero fuera del período analizado porque la empresa cierra el año por factura. Por alcance salen además las 3.090 ventas nacionales de 2025. El universo de la métrica son \textbf{10.672 líneas}.

\paragraph{Qué se corrigió sin excluir.} Dentro de ese universo la cadena rectificó 555 montos netos discordantes, 151 líneas cargadas en botellas en vez de cajas, 542 monedas vacías imputadas por mercado y 1.852 costos ausentes imputados por tier y añada (sobre el total registrado, las monedas vacías son 698). Ninguna corrección descarta la transacción: la conserva con la regla aplicada registrada en \texttt{output/cobertura.md}.

\paragraph{Supuestos.} Los trece supuestos que afectan alguna cifra están en el anexo de bitácora, con su justificación y su efecto si fueran falsos. Los que más pesan, S1, S3 y S8, están cuantificados y ninguno invierte la conclusión.
